% chapters/02/tujuan.tex
\begin{tujuanbox}
Bab ini tidak mengajarkan penulisan kode, melainkan membekali Anda kerangka
berpikir dan kosakata yang dipakai pada seluruh bab berikutnya. Setelah
menuntaskan bab ini, Anda mampu:
\begin{enumerate}
  \item menjelaskan permainan digital sebagai sistem waktu nyata beserta
        subsistem pembentuknya;
  \item menjelaskan hubungan antara bingkai, selang waktu, dan langkah
        simulasi, serta membedakan langkah tetap dan langkah berubah;
  \item menjelaskan sistem koordinat, skala, dan besaran turunan yang dipakai
        untuk menyatakan keadaan permainan;
  \item menjelaskan mengapa simulasi fisika pada komputer bersifat diskret dan
        bagaimana galat integrasi bergantung pada panjang langkah;
  \item membedakan lapisan keputusan dan lapisan navigasi pada pengendalian
        NPC;
  \item membedakan cara kerja pustaka dan mesin permainan, serta menjelaskan
        padanan konsep di antara keduanya.
\end{enumerate}
\end{tujuanbox}
